\section{Related Work}
\label{sec:related}

\paragraph{Common-variant GWAS of ASD.}
Grove et al.~\cite{grove2019} reported the first robust set of common-variant
risk loci for autism by combining 18\,381 cases and 27\,969 controls across
the iPSYCH and PGC cohorts (the dataset re-used here as the discovery
sumstats). Their PRS explained roughly 2.5\,\% of variance on liability
scale, leaving most ASD heritability unexplained -- evidence that more
sophisticated modelling, not just larger samples, is required.

\paragraph{Convergent biology from rare variants.}
Satterstrom et al.~\cite{satterstrom2020} sequenced exomes from 35\,584
individuals and identified 102 high-confidence ASD genes that cluster
around prenatal cortical development and synaptic
communication. Critically, these genes overlap the regions highlighted by
common-variant GWAS, suggesting that distinct sources of genetic risk
converge on a small set of biological systems -- a finding that motivates
gene-network-aware PRS methods.

\paragraph{Heritability gap.}
The Havdahl et al.~\cite{havdahl2021} review surveys SNP-based
heritability estimates ranging 12--65\,\% across studies, far below the
80\,\% from twin studies, and argues that closing this gap requires
models that integrate variant-level statistics with multi-omic context
rather than additive linear summation.

\paragraph{ML on focused SNP panels.}
Jiao et al.~\cite{jiao2012} demonstrated that decision-tree-based
classifiers trained on 29 SNPs from 9 candidate genes could predict
symptom severity (CARS) with 67\,\% accuracy in 118 children, establishing
that SNP patterns carry genuine ASD signal beyond simple presence/absence
of disease. The scope was, however, restricted to a tiny pre-selected
panel; whether genome-wide approaches can extend that finding remains
open.

\paragraph{ML-augmented polygenic scoring.}
Dong et al.~\cite{dong2023} combined PRS-CS-derived polygenic scores with
rare de novo mutation tallies in 9\,989 individuals and showed that the
two genetic components contribute through overlapping but distinct
pathways. Their pipeline kept the full SNP set rather than a thresholded
subset and made explicit that biology-aware modelling on summary
statistics is feasible.

\paragraph{Bayesian shrinkage PRS.}
PRS-CS~\cite{ge2019prscs} and LDpred2~\cite{prive2020ldpred2} are the
leading sumstats-only PRS construction methods. Both apply continuous or
discrete-mixture shrinkage to the marginal effect-size estimates while
respecting LD via reference-panel correlation matrices. Both, however,
treat SNPs as exchangeable conditional on LD, ignoring functional
annotations and gene-network structure.

\paragraph{Functional and network-aware extensions.}
LDSC partitioned heritability~\cite{finucane2015} and stratified-LD score
regression demonstrate that brain-specific functional categories carry
disproportionate heritability for psychiatric traits, motivating
annotation-aware PRS. MAGMA~\cite{deleeuw2015magma} provides a
gene-based test that aggregates SNP signals through a tissue-aware LD
projection. Graph-neural-network methods on protein--protein interaction
networks have been explored for cancer and Alzheimer's outcome
prediction, but to our knowledge no published ASD PRS uses a sumstats-only
GNN with the formal validation strategy described here.

\medskip\noindent
The methods we benchmark in Section~\ref{sec:methods} are positioned
directly against these works: P+T (the field standard threshold-based
construction) and a PRS-CS-style Bayesian baseline are compared against
two ML-enhanced alternatives that incorporate either SNP-level
annotations or explicit gene-network structure.
